\documentclass[pdflatex,sn-mathphys-num,draft]{sn-jnl}

\usepackage{graphicx}%
\usepackage{multirow}%
\usepackage{amsmath,amssymb,amsfonts}%
\usepackage{amsthm}%
\usepackage{mathrsfs}%
\usepackage[title]{appendix}%
\usepackage{xcolor}%
\usepackage{textcomp}%
\usepackage{manyfoot}%
\usepackage{booktabs}%
\usepackage{algorithm}%
\usepackage{algorithmicx}%
\usepackage{algpseudocode}%
\usepackage{listings}%
\usepackage{lmodern}%
\usepackage{siunitx}
\usepackage{microtype}
\usepackage{derivative}
\usepackage[version=4]{mhchem}
\usepackage{cleveref}

\renewcommand{\arraystretch}{1.2}
\sisetup{group-digits=true,
  group-separator={\,},
  separate-uncertainty
}

\theoremstyle{thmstyleone}%
\newtheorem{theorem}{Teorema}%
\newtheorem{proposition}[theorem]{Proposición}%

\theoremstyle{thmstyletwo}%
\newtheorem{example}{Ejemplo}%
\newtheorem{remark}{Observación}%

\theoremstyle{thmstylethree}%
\newtheorem{definition}{Definición}%

\DeclareSIUnit\angstrom{\text{\r{A}}}

\raggedbottom
%%\unnumbered% uncomment this for unnumbered level heads

\begin{document}

\title[Identificación y Caracterización del Nitrato de Plata]{
  Identificación y Caracterización Fisicoquímica del Nitrato de Plata: Una
  Revisión
}

\author{\fnm{Julian L.}
\sur{Avila-Martinez}}\email{jlavilam@udistrital.edu.co}

\affil{\orgdiv{Programa Académico de Física}, \orgname{Universidad Distrital
Francisco José de Caldas}, \orgaddress{\city{Bogotá D.C.},
\country{Colombia}}}

\abstract{
Presentamos una caracterización y revisión del nitrato de plata
($\ce{AgNO3}$), originada a partir de la identificación de un polvo
microcristalino blanco. El análisis experimental, mediante espectroscopia de
fluorescencia de rayos X (FRX), perfiles de solubilidad diferencial y pruebas
puntuales colorimétricas de aniones, identificó el material de manera
concluyente. Más allá de la identificación básica, el estudio proporciona un
examen detallado de la estructura de estado sólido del compuesto,
específicamente su cristalización en el grupo espacial ortorrómbico
centrosimétrico $Pbca$. Se presta especial atención a la desviación de los
modelos iónicos simples, impulsada por el poder polarizante del catión plata y
las interacciones argentofílicas débiles ($\ce{Ag \dots Ag}$). La discusión
fisicoquímica analiza su ruta de descomposición térmica hacia plata
elemental--distinta de los nitratos alcalinos--y la termodinámica de la depresión
del punto de fusión a nanoescala. Concluimos evaluando la versatilidad crítica
del compuesto en aplicaciones industriales, uniendo su papel histórico en la
fotografía analógica y la cauterización médica con su función contemporánea
como precursor primario para tintas conductoras en electrónica flexible.
}

\keywords{Nitrato de plata, análisis cualitativo inorgánico, fluorescencia de
rayos X}

\maketitle

\section{Introducción}
\label{sec:introduction}

El nitrato de plata (\ce{AgNO3}) ocupa una posición central en la química de
los metales de acuñación, sirviendo como el precursor preeminente para la
síntesis de prácticamente todos los demás derivados de la plata. Históricamente
denominado \textit{lapis infernalis} o ``cáustica lunar'' debido a su potente
capacidad cauterizadora y su asociación con el símbolo alquímico de la plata
(la luna) \cite{Szabadvary1992}, el compuesto ha trascendido sus orígenes
arcaicos para convertirse en una piedra angular de la química industrial y
analítica moderna \cite{Luceri2023, Sodhi2025, Alexander2009}.

Físicamente, el nitrato de plata es una sal inorgánica que desafía los modelos
iónicos simples. Si bien comparte la estequiometría de los nitratos de metales
alcalinos, el significativo poder polarizante del catión plata(I) ($\ce{Ag+}$)
introduce un marcado carácter covalente en su energía de red y en sus
interacciones de enlace. Esta estructura electrónica sustenta su fotosensibilidad
única--la propiedad fundamental de la industria fotográfica analógica--y su
extensa utilidad en electrónica, donde funciona como precursor de nanopartículas
de plata conductoras.

El alcance de este artículo sigue una progresión analítica deductiva. Comenzamos
con el examen de un polvo microcristalino blanco no identificado, utilizando
métodos cualitativos y espectroscópicos para establecer rigurosamente su
identidad química. Tras confirmar la sustancia como nitrato de plata, la
discusión se expande para proporcionar una revisión integral de sus propiedades
de estado sólido, comportamiento cristalográfico y su utilidad multifacética en
aplicaciones tecnológicas contemporáneas.

\section{Metodología Experimental para la Identificación}
\label{sec:methodology}

\subsection{Observación Física y Macroscópica}
La investigación comenzó con una muestra presentada como un polvo fino y blanco,
de naturaleza microcristalina. El material exhibió una marcada higroscopicidad,
formando aglomeraciones distintivas tras la exposición a la humedad atmosférica.
Las pruebas de susceptibilidad magnética, utilizando un imán de \ce{NdFeB} de alta
intensidad, no revelaron fuerzas atractivas o repulsivas observables,
caracterizando el material como no ferromagnético (diamagnético o débilmente
paramagnético).

Una observación diagnóstica ocurrió tras el contacto dérmico, donde el material
indujo una pigmentación oscura persistente en la epidermis. Este fenómeno es
consistente con la fotorreducción \textit{in situ} de iones de plata
($\ce{Ag+}$) a plata metálica elemental ($\ce{Ag^0}$), catalizada por reductores
orgánicos presentes en la piel. El análisis térmico demostró un punto de fusión
bajo sin evidencia de carbonización o combustión, clasificando la sustancia de
manera concluyente como una sal inorgánica en lugar de un compuesto orgánico.

\subsection{Composición Elemental y Perfil de Solubilidad}
El análisis elemental no destructivo se realizó mediante espectroscopia de
Fluorescencia de Rayos X (FRX) (ver \cref{fig:xrf_spectrum}). El espectro
resultante exhibió líneas de emisión prominentes características de la Plata
($\ce{Ag}$), identificándola como el constituyente catiónico primario.

\begin{figure}[htpb!]
  \centering
  \includegraphics[width=0.85\linewidth]{./images/xrf-spectrum.pdf}
  \caption{Espectro representativo de FRX de la muestra. El predominio de las
  líneas de emisión de plata (Series $K$ y $L$ de Ag) confirma la identidad
  catiónica, mientras que la ausencia de señales de halógenos pesados respalda
  la exclusión de haluros de plata insolubles.}
  \label{fig:xrf_spectrum}
\end{figure}

Las pruebas de solubilidad presentaron una aparente contradicción. El material
demostró alta solubilidad en agua desionizada ($\ce{H2O}$), formando una
solución clara e incolora. Por el contrario, la introducción en ácido sulfúrico
diluido ($\ce{H2SO4}$) resultó en la formación inmediata de un precipitado
sólido blanco. Esta observación no indica insolubilidad de la muestra original;
más bien, significa una reacción de precipitación que confirma la presencia de
iones $\ce{Ag+}$ libres:
\begin{equation}
  \ce{2Ag+(aq) + SO4^{2-}(aq) -> Ag2SO4(s)}
\end{equation}
La formación de sulfato de plata ($\ce{Ag2SO4}$), una sal poco soluble,
confirma que la muestra actúa como una fuente de iones de plata solubles.

\subsection{Determinación del Anión y Confirmación de Identidad}
Con el catión identificado como $\ce{Ag+}$ y la clase candidata restringida a
sales de plata inorgánicas solubles (principalmente $\ce{AgNO3}$, $\ce{AgF}$ o
$\ce{AgClO4}$), se empleó una prueba puntual específica para identificar el
anión. La aplicación del reactivo de difenilamina a la solución resultó en el
desarrollo rápido de una coloración azul intensa. Esta reacción de oxidación es
diagnóstica para la presencia del ion nitrato ($\ce{NO3-}$).

\subsection{Conclusión de Identidad}
La convergencia de la evidencia experimental--naturaleza higroscópica,
estabilidad térmica específica, confirmación del catión plata vía FRX,
precipitación característica de sulfato y una prueba positiva de difenilamina
para nitratos--identifica deductivamente el material desconocido como Nitrato de
Plata ($\ce{AgNO3}$).

\section{Análisis Cristalográfico y Estructural}
\label{sec:crystallography}

\subsection{Sistema Cristalino y Polimorfismo de Fase}
La estructura de estado sólido del nitrato de plata ($\ce{AgNO3}$) se desvía
significativamente de los modelos de empaquetamiento iónico simple
característicos de los nitratos de metales alcalinos \cite{Singh2017}. Esta
desviación surge del alto poder polarizante del catión plata(I) ($\ce{Ag+}$) y
la sustancial contribución covalente al enlace $\ce{Ag-O}$. El material exhibe
un polimorfismo dependiente de la temperatura, transitando de una fase
ortorrómbica ordenada en condiciones ambientales (Fase II) a una fase trigonal
desordenada a temperaturas elevadas (Fase I) \cite{osti_1279797}.

\subsection{Simetría del Grupo Espacial (Fase II)}
A temperatura y presión ambiente ($\approx \qty{298}{\kelvin}$), la red adopta
el sistema cristalino ortorrómbico. Si bien la literatura temprana referenciaba
ocasionalmente el grupo espacial acéntrico $P2_12_12_1$, la difracción de rayos
X de alta resolución y los cálculos teóricos han asignado definitivamente el
grupo espacial centrosimétrico $Pbca$ (N.º 61), con simetría de grupo puntual
$mmm$. Esta fase se define por una coordinación rígida y específica entre los
centros de plata y los ligandos de nitrato.

\subsection{Parámetros de Celda Unitaria}
La celda unitaria contiene $Z=8$ unidades formulares. Los parámetros de red para
el polimorfo $Pbca$ estándar son:
\begin{align}
  a &= \qty{6.992(2)}{\angstrom} \\
  b &= \qty{7.335(2)}{\angstrom} \\
  c &= \qty{10.125(2)}{\angstrom}
\end{align}
Estos valores reflejan la naturaleza anisotrópica del empaquetamiento
\cite{osti_1279797, Persson2012}.
Refinamientos computacionales alternativos arrojan ligeras variaciones ($a
\approx \qty{6.90}{\angstrom}$, $b \approx \qty{7.20}{\angstrom}$, $c \approx
\qty{10.01}{\angstrom}$), dependientes de la historia térmica específica de la
muestra cristalina.

\subsection{Geometría de Coordinación e Interacciones Electrónicas}
El ion $\ce{Ag+}$ funciona como un ácido de Lewis blando dentro de una red de
polímeros de coordinación. El centro de plata exhibe un número de coordinación
de 8, formando una geometría mejor descrita como un octaedro bicapado, donde las
caras de las tapas comparten una sola arista.

A diferencia de la alta simetría observada en el nitrato de sodio ($\ce{NaNO3}$),
las distancias de enlace $\ce{Ag-O}$ en $\ce{AgNO3}$ muestran una variación
significativa, oscilando entre \qty{2.37}{\angstrom} y \qty{2.60}{\angstrom}.
Esta distribución indica fuertes distorsiones tipo Jahn-Teller impulsadas por la
contribución covalente a la energía de red. El anión nitrato ($\ce{NO3-}$) actúa
como un ligando puente polidentado, conectando múltiples centros de plata en una
red tridimensional rígida.

Una característica crítica de la estructura electrónica es la presencia de
contactos cortos plata-plata. La distancia $\ce{Ag \dots Ag}$ es
aproximadamente \qty{3.227}{\angstrom}, cayendo dentro del rango de
interacciones argentofílicas débiles. Estas interacciones $d^{10}-d^{10}$
estabilizan la red ortorrómbica e influyen directamente en la estructura de
bandas electrónica del material, que se caracteriza por una brecha de banda
(\textit{band gap}) de aproximadamente \qty{2.01}{\electronvolt}.

\subsection{Transición a Alta Temperatura (Fase I)}
Al calentar a \qty{159.4}{\degreeCelsius} (\qty{432.6}{\kelvin}), el material
sufre una transición de fase reversible, displaziva y de primer orden. La
estructura se transforma a una simetría trigonal (romboédrica) con grupo
espacial $R3m$ (N.º 160). El impulsor principal de esta transición es el inicio
del desorden rotacional en los grupos nitrato; a medida que la energía térmica
supera la barrera rotacional, los iones $\ce{NO3-}$ planos exhiben libertad
orientacional, aumentando la entropía del sistema y su simetría efectiva.

\section{Propiedades Fisicoquímicas}
\label{sec:properties}

\subsection{Estabilidad Térmica y Cinética de Descomposición}
El perfil termodinámico del nitrato de plata se caracteriza por un punto de
fusión relativamente bajo y una vía de descomposición que diverge de la de los
nitratos de metales alcalinos. El compuesto se funde a \qty{212}{\degreeCelsius}
(\qty{485}{\kelvin}) para formar un líquido amarillento estable. Si bien la
descomposición inmediata es insignificante en el punto de fusión, una
disociación apreciable inicia a \qty{250}{\degreeCelsius} y procede hasta
completarse a \qty{444}{\degreeCelsius} (\qty{717}{\kelvin}).

A diferencia del nitrato de potasio, que se descompone a nitrito, el nitrato de
plata se descompone completamente al metal elemental. Esta inestabilidad surge
de la energía libre de formación positiva para el óxido de plata a temperaturas
elevadas. La reacción de descomposición es estequiométrica:
\begin{equation}
  \ce{2AgNO3(l) -> 2Ag(s) + 2NO2(g) + O2(g)}
\end{equation}
Esta reacción proporciona un método gravimétrico para la determinación de pureza
y actúa como el mecanismo fundamental para generar residuos de plata en el
ensayo al fuego (docimasia).

\subsection{Termodinámica de la Solvatación}
El nitrato de plata exhibe alta solubilidad en agua, una propiedad impulsada por
la alta entalpía de hidratación del catión plata que supera efectivamente la
energía de red. El proceso de disolución es endotérmico, indicando un mecanismo
impulsado por la entropía. La solubilidad aumenta marcadamente con la
temperatura, elevándose de \qty{122}{\gram}/\qty{100}{\milli\liter} a
\qty{0}{\degreeCelsius} hasta \qty{912}{\gram}/\qty{100}{\milli\liter} a
\qty{100}{\degreeCelsius}.

Las soluciones acuosas no son neutras; exhiben un rango de pH de
\numrange{5.4}{6.4} debido a la hidrólisis del ion plata:
\begin{equation}
  \ce{Ag+(aq) + H2O(l) <=> AgOH(aq) + H+(aq)}
\end{equation}
En medios no acuosos, la solubilidad varía significativamente. Aunque es soluble
en etanol (\qty{3.1}{\gram}/\qty{100}{\gram}), tales soluciones plantean
peligros de seguridad significativos debido a la potencial formación de
fulminatos explosivos \cite{CameoChemicals}. Por el contrario, la alta solubilidad en acetonitrilo
(\qty{111.8}{\gram}/\qty{100}{\gram}) lo convierte en un reactivo estándar para
electroquímica no acuosa.

\section{Aplicaciones Industriales y Tecnológicas}
\label{sec:applications}

\subsection{Fotografía y Sistemas de Imagen}
Históricamente la aplicación dominante del nitrato de plata, la industria
fotográfica depende de la reducción inducida por fotones de los haluros de
plata ($\ce{AgX}$). El nitrato de plata funciona como el precursor esencial para
la generación de estos haluros dentro de emulsiones de gelatina.

El mecanismo implica procesos mecánico-cuánticos: los fotones incidentes
interactúan con un grano de haluro, generando fotoelectrones que reducen los
iones de plata intersticiales a clústeres atómicos neutros. Estos clústeres
constituyen la ``imagen latente'', que cataliza la reducción de todo el grano
cristalino a plata metálica durante el revelado químico. Aunque los sensores
digitales han reemplazado a la película en los mercados de consumo, esta
química sigue siendo crítica para la imagen radiográfica de alta resolución y el
microfilm de archivo debido a su excepcional resolución espacial y estabilidad.

\subsection{Electrónica Impresa y Tintas Conductoras}
En el dominio de la electrónica flexible, el nitrato de plata sirve como materia
prima para la síntesis de nanopartículas de plata (AgNPs) y tintas reactivas.
Las formulaciones de tinta típicamente estabilizan los iones de plata con
alquilaminas \cite{Yang2017}; tras la impresión, el sustrato se somete a curado
térmico o químico.

La utilidad de estas tintas se basa en el fenómeno de depresión descrito por el
efecto Gibbs--Thomson \cite{Castro2005, Asoro2009}. La descomposición del complejo y
la posterior sinterización de las nanopartículas ocurren a temperaturas
compatibles con sustratos poliméricos ($< \qty{150}{\degreeCelsius}$), formando
una red de percolación conductora sin comprometer la integridad estructural del
dispositivo. Innovaciones recientes emplean vapor de agua para disminuir aún más
el umbral de sinterización, facilitando la fabricación de etiquetas RFID y
rejillas para celdas solares.

\subsection{Deposición Química y Catálisis}
La alta reflectividad de la plata elemental se explota en la fabricación de
espejos vía deposición no electrolítica (reacción de Tollens). El nitrato de
plata se compleja con amoníaco para prevenir la precipitación de óxido de plata,
luego se reduce mediante un aldehído (p. ej., glucosa o formaldehído) para
depositar una película metálica uniforme.

Más allá de la deposición, el nitrato de plata actúa como precursor para
catalizadores heterogéneos. Se utiliza para impregnar soportes de alúmina para
la epoxidación de etileno a óxido de etileno, un intermedio fundamental en la
producción de plásticos y anticongelantes.

\subsection{Utilidad Biológica y Médica}
Las aplicaciones médicas y biocidas del nitrato de plata derivan del efecto
oligodinámico, donde concentraciones traza de iones $\ce{Ag+}$ ejercen una
potente actividad antimicrobiana.

\begin{itemize}
  \item \textbf{Mecanismo Antimicrobiano:} Los iones de plata se unen con alta
    afinidad a grupos tiol ($\ce{-SH}$) en enzimas respiratorias bacterianas,
    interrumpiendo la respiración celular y generando especies reactivas de
    oxígeno (ROS) que inducen apoptosis \cite{Morones2005, Pandian2010}. Además,
    $\ce{Ag+}$ interactúa con el ADN bacteriano, condensando la hélice e
    inhibiendo la replicación.

  \item \textbf{Cauterización Terapéutica:} En forma sólida (lápiz de nitrato
    de plata), se emplea para la cauterización química para coagular vasos
    sanguíneos en epistaxis o para la ablación de tejido de hipergranulación en
    heridas en proceso de curación.

  \item \textbf{Integridad Electoral:} El nitrato de plata es el agente activo
    en la tinta indeleble electoral, utilizada para asegurar la integridad de
    los procesos de votación. Al contacto con la cutícula, reacciona con el
    cloruro de sodio y las proteínas de la piel para formar complejos
    fotosensibles. La exposición a la luz ultravioleta reduce estos complejos a
    plata metálica y sulfuro de plata, creando una pigmentación que persiste
    hasta la descamación epidérmica \cite{Nandwani2025}.
\end{itemize}

\section{Conclusión}
\label{sec:conclusion}

La caracterización sistemática presentada aquí identificó con éxito la sustancia
microcristalina desconocida como nitrato de plata (\ce{AgNO3}). Esta conclusión
se alcanzó a través de un marco analítico deductivo, donde la verificación
elemental vía FRX, las pruebas puntuales específicas de aniones y los
comportamientos únicos de solubilidad en ácido sulfúrico proporcionaron
evidencia convergente de la identidad del compuesto.

Más allá de la mera identificación, el análisis estructural resalta que el
nitrato de plata no es una sal iónica simple. Los datos cristalográficos
confirman una red ortorrómbica $Pbca$ centrosimétrica estabilizada por
contribuciones covalentes significativas \ce{Ag-O} e interacciones
argentofílicas débiles \ce{Ag \dots Ag}. Estas características estructurales
fundamentales gobiernan directamente su perfil fisicoquímico, más notablemente
su bajo punto de fusión e inestabilidad térmica, que lo distinguen marcadamente
de los nitratos de metales alcalinos.

En última instancia, este estudio subraya la versatilidad perdurable del nitrato
de plata. Aunque históricamente significativo como el fundamento de la
fotografía analógica y la cauterización médica, su capacidad única para
descomponerse en nanoestructuras metálicas conductoras a temperaturas moderadas
asegura su relevancia en la ciencia de materiales moderna. Desde la
esterilización oligodinámica de interfaces médicas hasta la fabricación de
circuitos electrónicos impresos, el nitrato de plata permanece como un precursor
indispensable en la síntesis de materiales funcionales avanzados.

\bibliography{./nano-particle-reference.bib}

\end{document}
